\section{Implementation Details}

\subsection{Model Architecture Details}

This appendix documents the implementation corresponding to the current MoM-Driver paper draft. The instantiated model is a camera-only multi-frame planning architecture composed of a DINOv2 visual encoder, a Mixture-of-Memories (MoM) temporal fusion backbone, and proposal-based trajectory decoding and scoring modules.

\paragraph{Inputs and tokenization.}
The model operates on a temporal context of $10$ frames collected from $4$ cameras, namely the front, back, left, and right views. Each image is resized to $574 \times 336$ and normalized using ImageNet statistics. The visual encoder is instantiated as \texttt{timm/vit\_small\_patch14\_reg4\_dinov2.lvd142m}, corresponding to DINOv2-Small with patch size $14$ and four register tokens.

For each camera stream, we prepend $16$ learnable scene tokens before the ViT encoder. Consequently, each frame yields $4 \times 16 = 64$ visual tokens, and the entire temporal window contains $640$ visual tokens before temporal fusion. The ego-state input is an $11$-dimensional vector comprising ego pose $(3)$, velocity $(3)$, acceleration $(3)$, and driving command $(2)$. In the default setting, only the latest ego-state vector is projected into the planning space, although the implementation also supports using the full temporal ego-state history.

\paragraph{Detailed layer configuration.}
Table~\ref{tab:mom_full_config} summarizes the default model configuration used in the current implementation.

\begin{table}[t]
\centering
\footnotesize
\setlength{\tabcolsep}{6pt}
\renewcommand{\arraystretch}{1}
\caption{Complete model configuration of MoM-Driver in the current implementation.}
\label{tab:mom_full_config}
\begin{tabularx}{\linewidth}{l l c X}
\toprule
\textbf{Module} & \textbf{Parameter} & \textbf{Value} & \textbf{Description} \\
\midrule

\multirow{7}{*}{\textbf{\shortstack[l]{Visual\\Encoder}}}
 & Visual encoder      & DINOv2-Small          & ViT-S/14, width $384$, 12 blocks \\
 & LoRA                & rank $32$             & Q/V adaptation in all blocks \\
 & Input resolution    & $574 \times 336$      & Per-camera image size \\
 & Cameras             & $4$                   & Front, back, left, right \\
 & Scene tokens        & $16$/cam              & Learnable camera tokens \\
 & Tokens per frame    & $64$                  & $4 \times 16$ scene tokens \\
 & Neck projection     & $384 \rightarrow 256$ & Visual-to-model projection \\
\midrule

\multirow{9}{*}{\textbf{\shortstack[l]{Temporal\\Encoder\\(MoM)}}}
 & Temporal window     & $10$                  & Historical frames \\
 & Temporal depth      & $2$                   & MoM temporal blocks \\
 & Hidden size         & $256$                 & Feature dimension \\
 & Attention heads     & $4$                   & Heads per block \\
 & Head dimension      & $64$                  & Dimension per head \\
 & Value expansion     & $1.0$                 & Value expansion ratio \\
 & Conv kernel         & $4$                   & Short convolution size \\
 & Output gate         & enabled               & Gated output \\
 & Positional embedding & $1\times64\times256$ & Temporal token embedding \\
\midrule

\multirow{4}{*}{\textbf{\shortstack[l]{Memory\\Config}}}
 & Memories            & $4$                   & Memory slots per block \\
 & Top-$k$ routing     & $2$                   & Selected memories per token \\
 & Shared projections  & enabled               & Shared memory branch \\
 & State-Frozen        & off                   & Implemented in code; enabled in the dedicated paper ablation \\
\midrule

\multirow{5}{*}{\textbf{\shortstack[l]{Trajectory\\Decoder}}}
 & Proposal tokens     & $64$                  & Learnable trajectory queries \\
 & Trajectory horizon  & $8$ poses             & $4$s horizon, $0.5$s interval \\
 & Trajectory decoder  & $4$ layers            & MoM decoder layers \\
 & Trajectory head     & $256 \!\rightarrow\! 1024 \!\rightarrow\! 1024 \!\rightarrow\! 24$ & Proposal prediction MLP \\
 & Trajectory heads    & $5$                   & Initial $+$ 4 refinements \\
\midrule

\multirow{4}{*}{\textbf{\shortstack[l]{Trajectory\\Scorer}}}
 & Scorer decoder      & $4$ layers            & MoM scorer layers \\
 & Scorer embedding    & $24 \!\rightarrow\! 1024 \!\rightarrow\! 256$ & Proposal re-embedding MLP \\
 & Score heads         & $6$                   & PDMS sub-score heads \\
 & Score head dim      & $256 \!\rightarrow\! 1024 \!\rightarrow\! 1$ & Per-score predictor \\

\bottomrule
\end{tabularx}
\end{table}

Each temporal MoM block applies RMSNorm before attention and before the gated MLP. The decoder layers adopt the same MoM configuration as the temporal backbone. Concretely, each decoder layer consists of (i) one MoM self-attention block operating over proposal tokens and (ii) one MoM cross-attention block that reads scene-conditioned temporal features. When State-Frozen cross-attention is enabled, query positions are prevented from writing into the recurrent state during cross-attention. This modification changes only the masking behavior inside the recurrent operator and does not introduce additional learnable parameters.

\paragraph{Parameterization.}
Under the default configuration, the full model contains approximately $57.46$M parameters, of which $35.55$M are trainable. The gap between total and trainable parameters arises because the DINOv2 backbone is frozen and only the LoRA adapters, scene tokens, temporal fusion layers, decoder layers, and prediction heads are optimized. The frozen DINOv2-Small backbone contributes approximately $21.91$M parameters, while the injected LoRA adapters contribute approximately $0.59$M trainable parameters. The two-layer temporal MoM backbone contains approximately $2.91$M parameters, each four-layer MoM decoder contributes approximately $11.65$M parameters, the trajectory heads contribute approximately $6.71$M parameters in total, and the planning score heads contribute approximately $1.59$M parameters.

\paragraph{Initialization strategy.}
The DINOv2 backbone is loaded from pretrained weights and remains frozen when LoRA is enabled. For each injected LoRA branch, the low-rank $A$ matrix is initialized with Kaiming uniform initialization, whereas the corresponding $B$ matrix is initialized to zero. This design guarantees that the initial LoRA-augmented model is functionally identical to the frozen pretrained backbone.

The learnable scene token bank is initialized from a zero-mean Gaussian distribution with standard deviation $10^{-6}$, and the temporal positional embedding is initialized from $\mathcal{N}(0, 0.02^2)$. Within each MoM attention layer, linear projections are initialized with Xavier uniform initialization using gain $2^{-2.5}$. The learned recurrent decay parameters $A_{\log}$ and $dt_{\mathrm{bias}}$ follow the default initialization scheme implemented in the FLA MoM operator. The final backbone LayerNorm is initialized with zero bias and unit scale. All remaining linear and embedding layers, including the ego-state encoder, proposal embeddings, scorer positional embedding, trajectory heads, and score heads, use the default PyTorch initialization.

\subsection{Training Details}

\paragraph{Data preprocessing and augmentation.}
Training is performed on cached NAVSIM samples with a temporal history length of $10$ frames. Images are resized to $574 \times 336$, converted to floating-point values in $[0,1]$, and normalized by ImageNet mean and standard deviation. The current implementation employs a single explicit image augmentation, namely GridMask, which is applied independently to each camera image during training. The GridMask configuration is as follows: horizontal masking enabled, vertical masking enabled, rotation $=1$, ratio $=0.5$, offset noise disabled, mode $=1$, and application probability $=0.7$.

No additional photometric distortion, color jitter, random crop, horizontal flip, MixUp, or CutMix augmentation is used in the current implementation. For variable-length histories, the input sequence is left-padded to length $10$, and the valid-frame count is propagated to the temporal fusion and decoder modules in order to mask padded frame tokens.

\paragraph{Optimizer and learning rate schedule.}
The implementation is optimized with AdamW using a base learning rate of $1 \times 10^{-4}$. In the current codebase, no explicit learning rate scheduler, warm-up phase, cosine annealing, or step decay is applied; the learning rate therefore remains constant throughout optimization. The default training configuration uses mixed precision (\texttt{16-mixed}), distributed data-parallel training, a maximum of $100$ epochs, no gradient clipping, and random seed $0$. The provided training script further overrides the batch size to $1$ and sets the number of dataloader workers to $16$.

For clarity, this implementation-level configuration is slightly more specific than the condensed training description in the main paper, which reports the paper-level setting as training for $100$ epochs with learning rate $1 \times 10^{-4}$ on $4\times$ RTX 4090 GPUs. The appendix numbers should therefore be read as the exact code configuration rather than a separate algorithmic change.

\paragraph{Loss formulation.}
The model predicts $64$ trajectory proposals and refines them over $4$ decoder stages. Let $\{\mathbf{P}^{(0)}, \dots, \mathbf{P}^{(R)}\}$ denote the initial proposal set and the $R=4$ refined proposal sets. At each stage, the trajectory loss is defined as the minimum proposal-to-ground-truth $L_1$ distance across the proposal dimension. When a long-horizon target is available, the implementation adds an additional minimum-distance term with respect to the interpolated long trajectory target.

The scorer predicts six planning-related attributes for each proposal: no-at-fault collision, drivable-area compliance, time-to-collision within bound, ego progress, driving-direction compliance, and comfort. These attributes are supervised using binary cross-entropy losses, with dedicated handling of ternary labels in the cached PDM supervision. The final planning score loss is the sum of the six attribute-level losses.

The overall training loss in the implementation is
\begin{equation}
\mathcal{L} =
\lambda_{\text{traj}} \mathcal{L}_{\text{traj}} +
\lambda_{\text{score}} \mathcal{L}_{\text{score}} +
\lambda_{\text{pred-ce}} \mathcal{L}_{\text{pred-ce}} +
\lambda_{\text{pred-l1}} \mathcal{L}_{\text{pred-l1}} +
\lambda_{\text{pred-area}} \mathcal{L}_{\text{pred-area}}.
\end{equation}
In the released default configuration, the auxiliary branches for area prediction, agent prediction, BEV prediction, and double-score prediction are disabled. As a result, only $\mathcal{L}_{\text{traj}}$ and $\mathcal{L}_{\text{score}}$ are active in practice, which is consistent with the simplified objective reported in the main paper.

\paragraph{Loss-weight tuning.}
The default loss weights follow the DrivoR-style setup adopted in the implementation:
\begin{equation}
\lambda_{\text{traj}} = 1.0,
\quad
\lambda_{\text{score}} = 1.0,
\quad
\lambda_{\text{inter}} = 0.0,
\quad
\lambda_{\text{prev}} = 0.0,
\quad
\lambda_{\text{pred-ce}} = 1.0,
\quad
\lambda_{\text{pred-l1}} = 0.1,
\quad
\lambda_{\text{pred-area}} = 2.0.
\end{equation}
However, because the corresponding auxiliary branches are disabled by default, the effective optimization objective reduces to trajectory regression and final planning-score supervision. A subtle but important implementation detail is that setting $\lambda_{\text{prev}} = 0$ disables recursive loss accumulation across refinement stages; consequently, the effective trajectory supervision is dominated by the final refinement stage.

For proposal ranking, the predicted proposal score is computed from the six attribute logits using the fixed aggregation weights
\begin{equation}
w_{\text{NAC}} = 1.0, \quad
w_{\text{DAC}} = 1.0, \quad
w_{\text{DDC}} = 0.0, \quad
w_{\text{TTC}} = 5.0, \quad
w_{\text{EP}} = 5.0, \quad
w_{\text{Comfort}} = 2.0.
\end{equation}
These coefficients are used for proposal ranking and score supervision and remain fixed throughout training.
